\section{语言与 LLM 驱动的高层调度}

\subsection{LLM 作为规划器的两个角色}

LLM 既可以输出子任务序列，也可以做 feasibility filter。讲者强调要把 LLM 生成的 plan 打上 provenance（来源时间、prompt hash），并对每个 step 进行 tool-check。

\begin{table}[H]
\centering
\begin{tabular}{p{4cm}p{5cm}}
\toprule
角色 & 实践要点 \\
\midrule
Plan generator & 打上 prompt id、temperature、llm version \\
Verifier & 依靠 tool chain 验证 plan 的可行性（如 python sandbox） \\
Skill mapper & 把 plan 转为 skill id 或 primitive action \\
\bottomrule
\end{tabular}
\caption{LLM 在层次化流程中的多重身份。}
\end{table}

\begin{importantbox}{LLM Prompt log}
每条 prompt 对应 metadata：`timestamp`, `llm-version`, `temperature`, `tool-call`, `expected skill id`，并将 hallucination/feasibility 反馈回 prompt log。
\end{importantbox}

\subsection{Plan verification pipeline}

\begin{figure}[H]
\centering
\includegraphics[width=0.85\textwidth]{slides-images/slide-33.jpg}
\caption{LLM 计划校验与 tool-check 流程（slide-33）。}
\end{figure}

每条 LLM 输出先进入 verifier（python sandbox + knowledge graph），再由 skill mapper 解析成可执行 plan。若发现 hallucination，立刻通知 incident channel 并把 context 记录入 grade-book。

\begin{warningbox}{Plan verification 的 guardrail}
\begin{itemize}
    \item 每一条 plan 包含 \texttt{plan\_id}、\texttt{llm\_version} 与 \texttt{prompt\_hash}；
    \item 执行前使用 tool-check 验证可行性；
    \item 发现 hallucination 时触发 rollback + 补 prompt 改动。
\end{itemize}
\end{warningbox}

\subsection{治理与审计}

生成规划时要有治理 guardrail，例如 multi-round verification、tool sandbox、human-in-the-loop review。若 LLM output 含 high-risk action（如 modify-files），必须触发 incident drill。

\begin{warningbox}{LLM 计划调度的风险}
自动化规划可能产生 hallucination，尤其在多步骤任务中。建议在部署前设定：
\begin{itemize}
    \item Multi-modal verification：用 tool call verify each step；
    \item Stop tokens：防止 LLM 无停止条件地追加 subtask；
    \item Human review：critical plan 前人工确认。
\end{itemize}
\end{warningbox}

\subsection{本章小结}

语言+LLM 让层次 RL 更具 explainability，但也把 hallucination/incident 的治理压力前置到高层 prompt log 和 verification pipeline。

